\section{The toric filling at the cusp}

\subsection{The infinite toric threefold}
Let
\[
 N'=\Z^2\oplus\Z,
 \qquad M'=\Hom(N',\Z),
\]
and write points of $N'_\R$ as $(y,y_3)$.  Let $\mathcal T_{A_2}$ be
the triangulation of $\R^2$ with vertices $\Z^2$ and triangles
\[
 \{v,v+e_1,v+e_2\},
 \qquad
 \{v+e_1,v+e_2,v+e_1+e_2\}.
\]
For every cell $P$ let $\wh P$ be the cone over $P\times\{1\}$, and let
$\Sigma$ be the fan consisting of these cones and their faces.

\begin{lemma}[The toric model]\label{lem:toric-model}
The collection $\Sigma$ is a smooth countable fan with support
$(\R^2\times\R_{>0})\cup\{0\}$.  The associated toric space
$Y=Y_\Sigma$ is a connected Hausdorff second-countable complex
threefold.  On its dense torus $(\C^*)^3$, use coordinates
$(x_1,x_2,t_c)$, where $t_c$ is the character corresponding to
$(0,0,1)\in M'$.  If a maximal cone has vertices $v_0,v_1,v_2$ and
$m_0,m_1,m_2$ is the dual basis, then
\[
 U_\sigma=\Spec\C[\wh\sigma^\vee\cap M']\cong\C^3,
 \qquad z_i=\chi^{m_i},
\]
and
\begin{equation}\label{eq:toric-local-product}
 t_c=z_0z_1z_2.
\end{equation}
The divisor $Y_0=t_c^{-1}(0)$ is reduced and equals
$\bigcup_{v\in\Z^2}E_v$.  Each $E_v$ is the toric del Pezzo surface
$dP_6$, and its toric boundary is the hexagon of six $(-1)$-curves indexed
by the six edges through $v$.
\end{lemma}

\begin{proof}
For the lower triangle the determinant of the three primitive generators
$(v,1),(v+e_1,1),(v+e_2,1)$ is $1$; for the upper triangle it is $-1$.
Thus every cone is unimodular.  Two cones meet in the cone over the common
face of their cells, and the support is as stated because a point
$(y,y_3)$ with $y_3>0$ lies over the cell containing $y/y_3$.
Glue the affine toric charts $U_\sigma$ along the open subspaces
$U_{\sigma\cap\tau}$.  The fan condition implies that if two points
represented in $U_\sigma$ and $U_\tau$ are identified through a chain of
charts, then they are already identified in $U_{\sigma\cap\tau}$.
Consequently the union of any two charts is the toric space of the finite
fan consisting of the faces of $\sigma$ and $\tau$; it is Hausdorff by
the standard separation argument for fans \cite[Secs.~1.4,3.1]{Fulton}.
Thus the full gluing is Hausdorff.  Every chart is smooth, the atlas is
countable, and all charts meet the dense torus.  Hence $Y$ is a connected
second-countable complex threefold.

The height character $(0,0,1)$ is nonnegative on the fan and takes value
one on every ray $(v,1)$.  Hence its divisor is
$\sum_vE_v$ with multiplicity one.  In a maximal chart,
$(0,0,1)=m_0+m_1+m_2$, proving
\eqref{eq:toric-local-product} and reducedness.  The star of the ray
$(v,1)$, modulo that ray, is the complete smooth fan with rays in the six
directions
$\pm e_1,\pm e_2,\pm(e_1-e_2)$, which is the fan of $dP_6$.
\end{proof}

\subsection{The lattice action}
Identify the cocharacter lattice of $(x_1,x_2,1)$ with
$\Lambda_{\mathrm{tor}}$ by $e_1\leftrightarrow\wh w$ and
$e_2\leftrightarrow\wh\delta$.  For
$\ol\lambda\in\ol\Lambda$ define
\[
 \varphi_{\ol\lambda}(y,y_3)
 =(y+y_3B_0\ol\lambda,y_3).
\]
It preserves the fan and induces a toric automorphism $\Phi_{\ol\lambda}$
of $Y$.  On the dense torus,
\[
 \Phi_{\ol\lambda}(x,t_c)
 =(t_c^{B_0\ol\lambda}x,t_c).
\]
Using \eqref{eq:cusp-exponential}, set
\begin{equation}\label{eq:Psi-action}
 \Psi_{\ol\lambda}(x,t_c)
 =\bigl(c_{\ol\lambda}(t_c)t_c^{B_0\ol\lambda}x,t_c\bigr).
\end{equation}
The formula extends holomorphically to $Y$ because it is the product of
$\Phi_{\ol\lambda}$ with a holomorphic torus translation.  The shears
$\varphi_{\ol\lambda}$ fix the horizontal cocharacter lattice
$\Z^2\oplus0$, so the toric automorphisms commute with these translations.
Since $C(t_c)(\ol\lambda+\ol\mu)=C(t_c)\ol\lambda+C(t_c)\ol\mu$, one has
$\Psi_{\ol\lambda}\Psi_{\ol\mu}=\Psi_{\ol\lambda+\ol\mu}$.  Thus
\eqref{eq:Psi-action} defines an action of $\ol\Lambda\cong\Z^2$.
The compact horizontal torus
$T_f=\Lambda_{\mathrm{tor}}\otimes S^1$ commutes with this action and
therefore acts on the quotient constructed below, preserving $t_c$.

For a torus point with $0<|t_c|<1$ define
\[
 y(p)=\frac{\log|x(p)|}{\log|t_c(p)|}\in\R^2.
\]
Let $\theta_i^\sigma$ be the barycentric coordinates of a maximal
triangle $\sigma$, and put
\[
 \sigma(\eta)=\{y:\theta_i^\sigma(y)>-\eta,
 \ i=0,1,2\}.
\]
For $S\geq1$ write
$U_\sigma(S)=\{p\in U_\sigma:|z_i(p)|<S\}$.

\begin{lemma}[Toric chart estimates]\label{lem:chart-estimates}
For a torus point $p$,
\[
 p\in U_\sigma(S)
 \quad\Longleftrightarrow\quad
 y(p)\in\sigma\left(\frac{\log S}{|\log|t_c(p)||}\right).
\]
Every point of $|t_c|<1$ lies in the closed unit polydisc of some maximal
toric chart.  In particular the sets $U_\sigma(2)$ cover $Y_\eps$, and
every compact subset has a finite subcover.
\end{lemma}

\begin{proof}
In the coordinates of \cref{lem:toric-model},
\[
 \log|z_i(p)|=\log|t_c(p)|\,\theta_i^\sigma(y(p)),
\]
which gives the first assertion because $\log|t_c|<0$.  For a torus point,
choose a triangle containing $y(p)$.

Let now $p$ lie in the orbit associated with the cone over a cell $P$.
For $\bar m\in M'\cap\widehat P^{\perp}$, the assignment
$\bar m\mapsto\log|\chi^{\bar m}(p)|$ is a real linear functional; under
the natural pairing it is represented by a vector
$\operatorname{Log}_P(p)$ in
$N'_\R/\operatorname{span}_\R\widehat P$.  For a maximal triangle
$\sigma\supset P$, let $\bar\sigma$ be the image of its cone in this
quotient.  The cones $\bar\sigma$, as $\sigma$ ranges over the maximal
triangles containing $P$, cover the quotient vector space: if $v$ is a
quotient vector and $q_0$ lies in the relative interior of
$P\times\{1\}$, then for sufficiently small $u>0$ the height-one
normalization of $q_0+u\widetilde v$ lies in a triangle containing $P$.
Choose $\sigma$ with
$-\operatorname{Log}_P(p)\in\bar\sigma$.  If $\bar m$ is a dual
coordinate on the orbit, then
$\log|\chi^{\bar m}(p)|=\langle\bar m,\operatorname{Log}_P(p)\rangle$;
the chosen inclusion therefore makes every nonzero chart coordinate have
modulus at most one.  This proves the second assertion.  The last assertion
follows from compactness.
\end{proof}

\begin{theorem}[Cusp filling]\label{thm:cusp-filling}
For sufficiently small $\eps>0$, the action of $\ol\Lambda$ on
\[
 Y_\eps=\{|t_c|<\eps\}\subset Y
\]
is free and properly discontinuous.  The quotient
\[
 N_0=Y_\eps/\ol\Lambda
\]
is a connected Hausdorff complex threefold, and $t_c$ descends to a proper
surjective holomorphic map
\[
 f_0:N_0\longrightarrow\Delta_0=\{|t_c|<\eps\}.
\]
It is a submersion away from $0$, and its nonzero fibres are the tori of
\cref{cor:smooth-family}.  The central fibre
\[
 W=f_0^{-1}(0)
\]
is reduced, irreducible and has normal crossings.  Its normalization is
$dP_6$; the normalization map identifies opposite sides of the
anticanonical hexagon.  The double locus consists of three rational curves
through two triple points.

Moreover
\[
 \pi_1(N_0)\cong\ol\Lambda,
\]
the inclusion of a nearby fibre sends $\Lambda$ to
$\ol\Lambda=\Lambda/\Lambda_{\mathrm{tor}}$, and the loop obtained by
holding $(x_1,x_2)=(1,1)$ and circling $t_c=0$ is null-homotopic in $N_0$.
\end{theorem}

\begin{proof}
Put
\[
 R(t_c)=-2\pi\Im C(t_c),
 \qquad
 B_{t_c}=B_0+\frac{R(t_c)}{\log|t_c|}.
\]
Choose $\eps$ so small that
$\|B_{t_c}-B_0\|\leq\frac12$ for $0<|t_c|<\eps$.
Since $B_0$ is orthogonal in the chosen bases,
\begin{equation}\label{eq:Bt-bound}
 |B_{t_c}\ol\lambda|\geq\frac12|\ol\lambda|.
\end{equation}

\smallskip
\noindent\emph{Freeness.}
On the dense torus, a fixed point of $\Psi_{\ol\lambda}$ with
$\ol\lambda\ne0$ would satisfy
\[
 \log|t_c|\,B_0\ol\lambda+R(t_c)\ol\lambda=0,
\]
which contradicts the choice of $\eps$.  On $Y_0$, the action sends the
orbit attached to a cell $P$ to the orbit attached to
$P+B_0\ol\lambda$.  No nonzero translation preserves a bounded cell, so
there are no fixed points there either.

\smallskip
\noindent\emph{Proper discontinuity.}
Equation \eqref{eq:Psi-action} gives, on the dense torus,
\begin{equation}\label{eq:position-translation}
 y(\Psi_{\ol\lambda}p)=y(p)+B_{t_c(p)}\ol\lambda.
\end{equation}
Fix maximal triangles $\sigma,\sigma'$.  By
\cref{lem:chart-estimates}, the positions of torus points in
$U_\sigma(2)\cap Y_\eps$ and $U_{\sigma'}(2)\cap Y_\eps$ lie in
bounded sets $\sigma(\eta_0)$ and $\sigma'(\eta_0)$, where
$\eta_0=\log2/|\log\eps|$.  If
$\Psi_{\ol\lambda}U_\sigma(2)$ meets $U_{\sigma'}(2)$, density of
the torus in every chart supplies a torus point in the intersection.
Then
\[
 \frac12|\ol\lambda|
 \leq |B_{t_c}\ol\lambda|
 \leq
 \sup\{|a-b|:a\in\sigma'(\eta_0),\ b\in\sigma(\eta_0)\},
\]
so only finitely many $\ol\lambda$ are possible.  A compact subset is
covered by finitely many $U_\sigma(2)$; therefore only finitely many of
its translates meet it.  This is proper discontinuity.

The orbit space is Hausdorff and the orbit map is a covering.  If $K$ is
a compact neighbourhood of a point $p$, proper discontinuity leaves only
finitely many nontrivial translates of $K$ that meet $K$; freeness and
Hausdorffness allow $K$ to be shrunk to an open neighbourhood disjoint
from all its nontrivial translates.  The same finite-translate argument
separates two distinct orbits.  Since all deck transformations are
biholomorphic, the quotient is a complex manifold.  The set $Y_\eps$ is connected: in the coordinates of a maximal chart,
$U_\sigma\cap Y_\eps=\{|z_0z_1z_2|<\eps\}$ is star-shaped and meets
the connected dense set
$(\C^*)^2\times\{0<|t_c|<\eps\}$.  It is second countable, and the open
quotient map shows that $N_0$ is second countable as well.  Hence $N_0$ is
a connected complex threefold.

\smallskip
\noindent\emph{Properness.}
Fix $0<\eps_1<\eps$.  The matrices $B_{t_c}$ and their inverses are
uniformly bounded for $0<|t_c|\leq\eps_1$.  Given a torus point $p$,
put $u=B_{t_c}^{-1}y(p)$ and choose $\ol\lambda\in\Z^2$ with
$|u+\ol\lambda|_\infty\leq1/2$.  Then
\[
 |y(\Psi_{\ol\lambda}p)|
 =|B_{t_c}(u+\ol\lambda)|
 \leq C
\]
for a constant independent of $p$.  Choose the finite set of maximal
triangles meeting the closed ball $\overline{B}(0,C+2)$ and, in each
corresponding chart, the compact unit polydisc
$\{|z_i|\leq1,\ |t_c|\leq\eps_1\}$.  By
\cref{lem:chart-estimates}, every translated torus point above the closed
disc lies in their union.

For a point of $Y_0$ in the orbit attached to a cell $P$, choose a vertex
$v$ of $P$.  Unimodularity of $B_0$ gives $\ol\lambda$ with
$B_0\ol\lambda=-v$, so the translated cell contains $0$ and is a face
of one of the finitely many triangles adjacent to $0$.  Enlarging the
preceding compact union by the corresponding closed chart polydiscs gives
a compact set $K_0$ meeting every orbit over $|t_c|\leq\eps_1$.
Therefore
$f_0^{-1}(|t_c|\leq\eps_1)$ is the compact image of $K_0$.  Any compact
subset of the open disc lies in such a closed subdisc, proving that $f_0$
is proper.

\smallskip
\noindent\emph{The fibres and local structure.}
For $t_c\ne0$, choose $s$ with $e^{2\pi is}=t_c$.  The exponential map
$\C^2\to(\C^*)^2$ identifies the quotient fibre with
\[
 \C^2/\bigl(\Z^2+(sB_0+C(t_c))\Z^2\bigr)
 =\C^2/\Pi(z)\Lambda,
\]
where $s=s(z)$; this is the torus family of
\cref{cor:smooth-family}.  These fibres exist
for every $0<|t_c|<\eps$, while $Y_0$ is nonempty, so $f_0$ is
surjective.  On $t_c\ne0$ the character $t_c:(\C^*)^3\to\C^*$ is a
submersion, and the quotient map is locally biholomorphic.  Hence $f_0$
is a submersion away from the origin.  Finally, suppose that exactly the coordinates indexed by a nonempty set
$I\subset\{0,1,2\}$ vanish at a chosen point.  The product of the remaining
coordinates is a holomorphic unit; absorbing that unit into one coordinate
indexed by $I$, \eqref{eq:toric-local-product} gives local coordinates on
$N_0$ in which $f_0$ is $z_0$, $z_0z_1$, or $z_0z_1z_2$, according as
$|I|=1,2,3$.  Thus $W$ is reduced and has normal crossings.

Because $B_0\ol\Lambda=\Z^2$, the action on the height-one lattice
is the full translation action.  Every cell has a vertex, so every toric
orbit in $Y_0$ is equivalent to an orbit contained in $E_0$; hence the map
$E_0\to W$ is surjective.  It is injective on the open torus orbit.  Two
orbits in $E_0$ can be identified only when cells $P,P'$ containing $0$
satisfy $P'=P+v$ for a nonzero integral translation.  For edges this
forces $P=\{0,v\}$ and $P'=\{0,-v\}$.  No nonzero translation preserves
an edge, so each boundary curve maps isomorphically to its image.  Thus
the six boundary curves are identified in the three opposite pairs and no
other curve identifications occur.  For triangles, the six triangles through $0$ fall into two
translation orbits, alternating around the hexagon; their images are the
two triple points.  The fibres of $E_0\to W$ are therefore finite.  The map is proper because
$E_0$ is compact and $W$ is Hausdorff; a proper holomorphic map with finite
fibres is finite.  Since $E_0$ is irreducible and maps onto $W$, the surface
$W$ is irreducible.  The finite map $E_0\to W$ is generically one-to-one,
with normal source and reduced target; it is therefore the normalization.  This proves the stated
description of the double and triple loci.

\smallskip
\noindent\emph{Fundamental group.}
The toric threefold $Y$ is simply connected.  For $n\ge1$, let
$\Sigma_n$ be the finite subfan consisting of the cones over the cells of
$\mathcal T_{A_2}$ contained in the square $[-n,n]^2$, together with all
their faces.  The corresponding toric spaces $Y_{\Sigma_n}$ form an
increasing open cover of $Y$.  Each contains the rays through
$(0,0,1),(1,0,1),(0,1,1)$, which generate $N'$ over $\Z$.  The standard
toric fundamental-group formula \cite[Sec.~3.2]{Fulton} therefore gives
$\pi_1(Y_{\Sigma_n})=0$.  Every loop has compact image and hence lies
in some $Y_{\Sigma_n}$, so $\pi_1(Y)=0$.

For $0<u\leq1$, the torus element $(1,1,u)\in(\C^*)^3$ acts on $Y$ and
satisfies $t_c((1,1,u)\cdot p)=u\,t_c(p)$.  A loop in $Y_\eps$ bounds a
disc in $Y$; for $u>0$ sufficiently small, translating that disc by
$(1,1,u)$ places it in $Y_\eps$.  On the boundary, the path of torus
elements $(1,1,u_s)$, with $u_s$ decreasing from $1$ to $u$, gives a
homotopy inside $Y_\eps$ from the original loop to the translated one.
Hence $Y_\eps$ is simply connected.
Therefore the covering $Y_\eps\to N_0$ is universal and
$\pi_1(N_0)=\ol\Lambda$.

On a nonzero fibre the covering $(\C^*)^2\to f_0^{-1}(t_c)$ has deck group
$\ol\Lambda$, so the induced map on fundamental groups is the projection
$\Lambda\to\ol\Lambda$.  Finally, the curve
$\{(1,1,t_c):|t_c|<\eps\}$ is a holomorphic disc in a toric chart; its
boundary is the stated meridian, which is therefore null-homotopic.
\end{proof}


\begin{proposition}[Identification over the punctured cusp disc]
\label{prop:cusp-identification}
After decreasing $\eps$ if necessary, put
$U_\eps=U_r\cap\pi^{-1}(D_0^*)$.  The map
\begin{equation}\label{eq:cusp-exp-map}
 E_0^{\exp}:U_\eps\times\C^2\longrightarrow
 Y_\eps\cap(\C^*)^3,
 \qquad
 E_0^{\exp}(z,\zeta)=
 \bigl(e^{2\pi i\zeta_1},e^{2\pi i\zeta_2},e^{2\pi is(z)}\bigr),
\end{equation}
induces a biholomorphism
\begin{equation}\label{eq:G0}
 G_0:J|_{D_0^*}\xrightarrow{\sim}N_0\setminus W,
 \qquad D_0^*=\{0<|t_c|<\eps\}.
\end{equation}
It identifies the marking $H_1(F;\Z)=\Lambda$ on both sides and carries the
zero section to the section represented by $(x_1,x_2)=(1,1)$.  The latter
extends across $t_c=0$ as a holomorphic section of $f_0$, meeting $W$ in
its smooth locus.
\end{proposition}

\begin{proof}
Let $\wt g_0(z,\zeta)=(g_0z,\zeta)$; the latter
formula follows from $R_0=I_2$.  Let $H$ be the subgroup generated by
$\wt g_0$ and the translations $t_\lambda$ with
$\lambda\in\Lambda_{\mathrm{tor}}$.  It is normal in
$\Lambda\rtimes\langle g_0\rangle$: conjugation by $\wt g_0$ preserves
the translations by $\Lambda_{\mathrm{tor}}$, while
\[
 t_\mu\,\wt g_0\,t_\mu^{-1}
 =t_{\mu-M_0\mu}\,\wt g_0\in H
 \qquad(\mu\in\Lambda),
\]
because $(M_0-I)\Lambda\subset\Lambda_{\mathrm{tor}}$.

The map $s:U_\eps\to\{\Im s>c\}$ is a biholomorphism and satisfies
$s\circ g_0=s-1$.  Hence the last coordinate in
\eqref{eq:cusp-exp-map} is the universal covering of the punctured disc,
with deck group generated by $g_0$.  The first two coordinates are the
ordinary exponential covering $\C^2\to(\C^*)^2$, whose deck lattice is
$\Z^2=\Pi(z)\Lambda_{\mathrm{tor}}$.  It follows that
$E_0^{\exp}$ is a surjective local biholomorphism and that its fibres are
exactly the $H$-orbits.

Write $\lambda=\lambda_{\mathrm{tor}}+\lambda'$ and let
$\ol\lambda$ be its image in $\ol\Lambda$.  By
\eqref{eq:cusp-normal-form} and \eqref{eq:cusp-exponential},
\[
 E_0^{\exp}\bigl(z,\zeta+\Pi(z)\lambda\bigr)
 =\Psi_{\ol\lambda}\bigl(E_0^{\exp}(z,\zeta)\bigr).
\]
Thus $E_0^{\exp}$ identifies $(U_\eps\times\C^2)/H$ with
$Y_\eps\cap(\C^*)^3$ and intertwines the residual action
$(\Lambda\rtimes\langle g_0\rangle)/H\cong\ol\Lambda$ with the action
$\Psi$ of \eqref{eq:Psi-action}.  Since $U_\eps$ is the only component of
$\pi^{-1}(D_0^*)$ stabilized by $g_0$ and its other $\Delta$-translates
are disjoint, the quotient by the full residual action is precisely
\eqref{eq:G0}.

The construction sends the cycles in $\Lambda_{\mathrm{tor}}$ to the
fundamental cycles of $(\C^*)^2$ and the quotient
$\Lambda/\Lambda_{\mathrm{tor}}$ to the deck group of a fibre of
$N_0\setminus W$; hence it preserves the marking.  Finally
$E_0^{\exp}(z,0)=(1,1,t_c)$.  In the toric chart corresponding to the
triangle with vertices $0,e_1,e_2$, the coordinates are
$(t_c/(x_1x_2),x_1,x_2)$, so
\[
 \{(1,1,t_c):|t_c|<\eps\}
\]
is a holomorphic disc meeting the divisor $t_c=0$ transversely at a point
of the open orbit of $E_0$.  Since every $\Psi_{\ol\lambda}$ preserves $t_c$ and the disc contains
exactly one point over each value of $t_c$, the quotient map is injective
on this disc.  Its image is therefore the required holomorphic section.
\end{proof}

\begin{proposition}[Boundary of the cusp piece]\label{prop:cusp-boundary}
For $0<\rho<\eps$, the boundary
$\mathcal M_0=f_0^{-1}(|t_c|=\rho)$ is the mapping torus of the monodromy
$M_0$ on $F$.  With $a_0$ the toric meridian,
\[
 \pi_1(\mathcal M_0)=\Lambda\rtimes_{M_0}\Z\langle a_0\rangle.
\]
The inclusion $\mathcal M_0\hookrightarrow N_0$ sends
$\lambda\mapsto\ol\lambda\in\ol\Lambda$ and $a_0\mapsto1$; its
kernel is the normal closure of
$\Lambda_{\mathrm{tor}}\cup\{a_0\}$.
\end{proposition}

\begin{proof}
Ehresmann's theorem identifies the restriction over the circle with the
mapping torus of the geometric monodromy, whose translation part is
isotopic to zero and whose linear part is $M_0$.  The map on the fibre
fundamental group is the projection established in
\cref{thm:cusp-filling}, and the meridian is null-homotopic there.  The
resulting quotient of the displayed semidirect product is
$\Lambda/\Lambda_{\mathrm{tor}}=\ol\Lambda$, which is
$\pi_1(N_0)$; hence the stated normal subgroup is exactly the kernel.
\end{proof}

